\documentclass[12pt]{article}
\usepackage{amsfonts,amsthm,amsmath,amssymb,mathtools}
\usepackage{mathrsfs}
\usepackage{enumitem}
\usepackage{tikz-cd}
\usepackage{geometry}
\usepackage{mlmodern}

\geometry{letterpaper, portrait, margin=1in}

\setcounter{section}{0}

\renewcommand{\thesection}{\S\arabic{section}}
\renewcommand{\thesubsection}{\arabic{section}}
\newcommand{\dd}{\mathop{}\!\mathrm{d}}
\newcommand{\R}{\mathbb{R}}
\newcommand{\Z}{\mathbb{Z}}
\newcommand{\C}{\mathbb{C}}
\newcommand{\Q}{\mathbb{Q}}

\DeclareMathOperator{\lcm}{lcm}
\DeclareMathOperator{\im}{im}

\newtheorem{thm}{Theorem}
\newenvironment{theorem}[1]{
	\renewcommand{\thethm}{#1}
	\begin{thm}
		}{\end{thm}}

\newtheorem{lma}{Lemma}
\newenvironment{lemma}[1]{
	\renewcommand{\thelma}{#1}
	\begin{lma}
		}{\end{lma}}

\theoremstyle{definition}
\newtheorem{ex}{}
\newenvironment{exercise}[1]{
	\renewcommand{\theex}{#1}
	\begin{ex}
		}{\end{ex}}

\title{Homework 11}
\author{Connor Mooneyhan}
\date{November 13, 2025}

\begin{document}

\begin{center}
	{\bf\Large MTH 346/646 homework cover sheet}
\end{center}

\vspace{0.5in}

\noindent{\large Name}: Connor Mooneyhan \\
{\large Homework number}: 11 \\
{\large Date}: November 13, 2025 \\

{\Large What topics did this assignment cover? How comfortable do you feel with
each of those topics?}


\vspace{1.75in}

{\Large In the course of doing this assignment, where did you get stuck?
	How did you get unstuck?}


\vspace{1.75in}

{\Large Is there anything you've done here that you're not confident about, that you want me to take a really close look at, or that you want me to suggest an alternate approach to?}


\pagebreak

\begin{ex}
	Let $C : y^{2} = x^{3} - 49042x$. Compute the rank of $C$ by computing
	the image of $\alpha$ and $\overline{\alpha}$ with assistance from Magma,
	but do not ask Magma for the rank or the Mordell-Weil group. (This is like
	the last example from the notes on 11/14.)
\end{ex}
\begin{proof}
	The prime factorization of $49042$ is $2 \cdot 7 \cdot 31 \cdot 113$, so the divisors are \begin{align*}
		D = \lbrace \pm 1, \pm 2, \pm 7, \pm 14, \pm 31, \pm 62, \pm 113, \pm 217, \pm 226, \pm 434, \\
		\pm 791, \pm 1582, \pm 3503, \pm 7006, \pm 24521, \pm 49042 \rbrace
	\end{align*}

	I used JavaScript to find the following points, given as solutions to the equation $N^2 = dM^4 - (49042/d) e^4$: \begin{equation*}
		\begin{array}{r|c}
			\textbf{Divisor} & \textbf{(M, e, N)}       \\
			\hline
			1                & (1, 0, 1)                \\
			-1               & (5289, 1999, 884489296)  \\
			2                & (11, 1, 69)              \\
			-2               & (13105, 1387, 178211938) \\
			7                & (7, 1, 99)               \\
			-7               & (2902, 2550, 543815019)  \\
			14               & (4, 1, 9)                \\
			-14              & (7619, 3081, 518145689)  \\
			31               & (86508, 4, 41667111046)  \\
			-31              & (1, 2, 159)              \\
			62               & (19798, 0, 3086302457)   \\
			-62              & (1, 1, 27)               \\
			113              & (3, 2, 47)               \\
			-113             & (4442, 4233, 308785241)  \\
			217              & (98892, 0, 144062911405) \\
			-217             & (1, 1, 3)                \\
			226              & (1, 1, 3)                \\
			-226             & (1279, 8908, 1168676075) \\
			434              & (89854, 0, 168197561503) \\
			-434             & (2, 3, 47)               \\
			791              & (26243, 0, 19369356948)  \\
			-791             & (3237, 6399, 130798511)  \\
			1582             & (45090, 0, 80865580813)  \\
			-1582            & (901, 10405, 601923169)  \\
			3503             & (65544, 0, 254264630822) \\
			-3503            & (1, 4, 9)                \\
			7006             & (71111, 0, 423261374961) \\
			-7006            & (1, 7, 99)               \\
			24521            & (62521, 0, 612097943287) \\
			-24521           & (1, 11, 69)              \\
			49042            & (37918, 0, 318401372375) \\
			-49042           & (2, 43, 1623).
		\end{array}
	\end{equation*}
	So in fact $\im\alpha = D$ and $|\im\alpha| = 32$.

	Now for $\overline{C}: y^2 = x^3 + 196168$, we consider only positive divisors $d$ which satisfy $N^2 = dM^4 + (196168/d)e^4$.
	This gives us the list \begin{equation*}
		\begin{array}{r|c}
			\textbf{Divisor} & \textbf{(M, e, N)}        \\
			\hline
			1                & (1, 0, 1)                 \\
			2                & (4, 1, 314)               \\
			4                & (1, 0, 2)                 \\
			7                & (78526, 11, 16314582762)  \\
			8                & (2, 1, 157)               \\
			14               & (81457, 5, 24826805419)   \\
			28               & (82577, 5, 36082549636)   \\
			31               & (96105, 13, 51424823882)  \\
			56               & (81457, 10, 49653610838)  \\
			62               & (19798, 0, 3086302457)    \\
			113              & (1, 1, 43)                \\
			124              & (84857, 9, 80183718050)   \\
			217              & (98892, 0, 144062911405)  \\
			226              & (6, 1, 542)               \\
			248              & (19798, 0, 6172604914)    \\
			434              & (89854, 0, 168197561503)  \\
			452              & (1, 2, 86)                \\
			791              & (26243, 0, 19369356948)   \\
			868              & (98892, 0, 288125822810)  \\
			904              & (56897, 0, 97333636911)   \\
			1582             & (45090, 0, 80865580813)   \\
			1736             & (89854, 0, 336395123006)  \\
			3164             & (26243, 0, 38738713896)   \\
			3503             & (65544, 0, 254264630822)  \\
			6328             & (45090, 0, 161731161626)  \\
			7006             & (71111, 0, 423261374961)  \\
			14012            & (32772, 0, 127132315411)  \\
			24521            & (62521, 0, 612097943287)  \\
			28024            & (62438, 0, 652623982229)  \\
			49042            & (37918, 0, 318401372375)  \\
			98084            & (34993, 0, 383496529291)  \\
			196168           & (37918, 0, 636802744750).
		\end{array}
	\end{equation*}
	This is also of size 32, so we should have \begin{equation*}
		2^r = \frac{|\im\alpha||\im\overline{\alpha}|}{4} = \frac{32\cdot 32}{4} = 2^8,
	\end{equation*}
	so we get the rank $r = 8$.
\end{proof}

\begin{ex}
	Use Magma to determine if there is a rational point on
	\[
		y^{2} = -x^{4} + 25268
	\]
	by asking Magma to compute the rank of an elliptic curve and thinking
	about the images of $\alpha$ and $\overline{\alpha}$ for that elliptic curve.
\end{ex}

\begin{ex}
	$*$ This exercise goes through the proof that there are no rational points
	on $2y^{2} = x^{4} - 17$, which shows that if $C : y^{2} = x^{3} + 17x$,
	then $2 \not\in \im \overline{\alpha}$. [I got this exercise from Nigel Boston.]

	\begin{enumerate}[label=(\alph*)]
		\item Suppose that $(a,b)$ is rational point on the curve $x^{4}
			      - 17 = 2y^{2}$. Show that $a = A/C$, $b = B/C^{2}$, where $A, B, C \in \Z$,
		      $\gcd(A,B) = \gcd(A,C) = \gcd(B,C) = 1$.

		\item Rewriting the equation as $(5A^{2} + 17C^{2})^{2} - (4B)^{2} =
			      17(A^{2} + 5C^{2})^{2}$, show that there are $U,V \in \Z$ such that
		      \[
			      5A^{2} + 17C^{2} \pm 4B = 17U^{2}, 5A^{2} + 17C^{2} \mp 4B = V^{2},
			      A^{2} + 5C^{2} = UV,
		      \]
		      or
		      \[
			      5A^{2} + 17C^{2} \pm 4B = 34U^{2}, 5A^{2} + 17C^{2} \mp 4B = 2V^{2},
			      A^{2} + 5C^{2} = 2UV.
		      \]

		\item Hence show that (eliminating $B$) either
		      \[
			      10A^{2} + 34C^{2} = 17U^{2} + V^{2}, A^{2} + 5C^{2} = UV
		      \]
		      or
		      \[
			      5A^{2} + 17C^{2} = 17U^{2} + V^{2}, A^{2} + 5C^{2} = 2UV.
		      \]

		\item Show that there are no solutions to $x^{2} \equiv 10 \pmod{17}$,
		      and no solutions to $x^{2} \equiv 5 \pmod{17}$ and use this to conclude
		      that in any solution to the equations in (c) one must have $17 | A$ and $17 | C$. This contradicts (a).
	\end{enumerate}
\end{ex}
\begin{proof}\
	\begin{enumerate}[label=(\alph*)]
		\item Here's what I've got.
		      Write $a = \alpha/\beta$ and $b = \gamma/\delta$ so that each fraction is reduced.
		      Then \begin{align*}
			               & \frac{\alpha^4}{\beta^4} - 17 = 2\cdot \frac{\gamma^2}{\delta^2} \\
			      \implies & \alpha^4\delta^4 - 17\beta^4\delta^4 = 2\gamma^2\delta^2,
		      \end{align*}
          so we can set $A = \alpha\delta$, $C = \beta\delta$, and $B = \gamma\delta$.
        \item 
	\end{enumerate}
\end{proof}

\end{document}
